\documentclass[a4paper, serif, 10pt]{moderncv}

%% moderncv
% color and style
\moderncvcolor{black}
\moderncvstyle{banking}

%% page layout/margins
\usepackage[top=2.5cm, bottom=2.5cm, left=1.5cm, right=1.5cm]{geometry}

\nopagenumbers{}

%% Babel config for Spanish language
% \usepackage[spanish]{babel} has some conflicting issues with moderncv
% so we have to use enable the following options to make it work
%
% ref:
% - https://tex.stackexchange.com/a/140161/36007
\usepackage[spanish,es-lcroman]{babel}

%% fontspec
\usepackage{fontspec}

\IfFontExistsTF{Linux Libertine}{
  \setmainfont[Ligatures={TeX, Common}, Numbers=OldStyle]{Linux Libertine}
}{}
\IfFontExistsTF{Linux Libertine O}{
  \setmainfont[Ligatures={TeX, Common}, Numbers=OldStyle]{Linux Libertine O}
}{}

%% CTeX
% CJK support, used to show CJK characters in the resume
%
% - heading=false: disable ctex heading
% - punct=kaiming: use kaiming punctuations styles for CJK
% - scheme=plain: use plain scheme, do not override \normalsize font size
% - space=auto: space settings for CJK characters
%
% CJK fonts are resolved per locale language, e.g. Noto Serif/Sans CJK JP
% for Japanese resumes, with fallback to the next region if not installed
%
% ref:
% - http://ctan.mirrorcatalogs.com/language/chinese/ctex/ctex.pdf
\usepackage[UTF8, heading=false, punct=kaiming, scheme=plain, space=auto]{ctex}

\IfFontExistsTF{Noto Serif CJK SC}{
  \setCJKmainfont{Noto Serif CJK SC}
}{}
\IfFontExistsTF{Noto Sans CJK SC}{
  \setCJKsansfont{Noto Sans CJK SC}
}{}

%% line spacing
\usepackage{setspace}
\setstretch{1.125}

%% URL styling - use normal text instead of monospace
\urlstyle{same}

%% Patch \httplink and \httpslink to handle full URLs
% The original moderncv macros always prepend http:// or https:// to URLs.
% This causes issues when the URL already has a protocol (e.g., https://example.com)
% resulting in malformed URLs like https://https://example.com.
% This patch checks if the URL already contains "://" and uses it directly if so.
% Note: We use \str_set:Nx (x-expansion) to expand macros like \@homepage first.
\ExplSyntaxOn
\str_new:N \g_yamlresume_url_str

\RenewDocumentCommand{\httpslink}{O{}m}{
  \str_set:Nx \g_yamlresume_url_str {#2}
  \str_if_in:NnTF \g_yamlresume_url_str {://}
    { \url{#2} }
    { \url{https://#2} }
}

\RenewDocumentCommand{\httplink}{O{}m}{
  \str_set:Nx \g_yamlresume_url_str {#2}
  \str_if_in:NnTF \g_yamlresume_url_str {://}
    { \url{#2} }
    { \url{http://#2} }
}
\ExplSyntaxOff

\name{Carlos Mendoza}{}
\title{Ingeniero de Ciberseguridad Senior}
\phone[mobile]{+34 612 345 678}
\email{carlos.mendoza@example.com}
\homepage{https://www.carlosmendoza.dev}

\address{Calle de Alcalá 123, Madrid, Comunidad de Madrid, España, 28009}{}{}

\extrainfo{{\small \faLinkedin}\ \href{https://www.linkedin.com/in/carlosmendoza}{@carlosmendoza} {} {} {} • {} {} {} 
{\small \faGithub}\ \href{https://github.com/cmendoza-sec}{@cmendoza-sec} {} {} {} • {} {} {} 
{\small \faTwitter}\ \href{https://twitter.com/cmendoza\_sec}{@cmendoza\_sec}}

\begin{document}

\maketitle

\section{Información básica}

\cvline{}{\begin{itemize}
\item Ingeniero de Ciberseguridad con más de 5 años de experiencia en la protección de infraestructuras críticas y datos sensibles.
\item Especializado en pentesting, análisis de vulnerabilidades y respuesta a incidentes, con certificaciones OSCP y CISSP.
\item Sólidos conocimientos en herramientas SIEM, firewalls y sistemas de detección de intrusiones.
\item Apasionado por la seguridad ofensiva y la formación de equipos, con habilidades de comunicación y trabajo en equipo.
\end{itemize}}
\section{Educación}

\cventry{sept 2015–jul 2019}
        {Licenciatura, Ingeniería Informática, Puntuación: 8.9}
        {Universidad Politécnica de Madrid}
        {\url{https://www.upm.es}}
        {}
        {\begin{itemize}
\item Desarrollé habilidades en programación (Python, C, Java) y análisis de sistemas.
\item Adquirí experiencia práctica en seguridad mediante proyectos de laboratorio y competiciones CTF.
\item Participé en proyectos grupales, mejorando mi comunicación y trabajo en equipo.
\end{itemize}
\textbf{Cursos}: Fundamentos de Seguridad Informática, Criptografía, Redes de Computadores, Sistemas Operativos, Programación en Python, Bases de Datos, Arquitectura de Computadores}
\section{Experiencia laboral}

\cventry{ene 2022 hasta la fecha}
        {Ingeniero de Ciberseguridad}
        {Telefónica Tech}
        {\url{https://telefonicatech.com}}
        {}
        {\begin{itemize}
\item Lideré pruebas de penetración y auditorías de seguridad en aplicaciones web y móviles.
\item Implementé y gestioné soluciones SIEM para la detección temprana de amenazas.
\item Coordiné la respuesta a incidentes de seguridad, reduciendo el tiempo de contención en un 30\%.
\item Desarrollé políticas de seguridad y procedimientos para el cumplimiento normativo (ISO 27001, GDPR).
\end{itemize}
\textbf{Palabras clave}: Pentesting, SIEM, Gestión de vulnerabilidades, Respuesta a incidentes, Cumplimiento normativo}

\cventry{sept 2019–dic 2021}
        {Analista de Seguridad}
        {Banco Santander}
        {\url{https://www.santander.com}}
        {}
        {\begin{itemize}
\item Realicé análisis de riesgos y evaluaciones de seguridad en infraestructuras críticas.
\item Configuré y mantuve firewalls, IDS/IPS y sistemas de prevención de intrusiones.
\item Colaboré en la elaboración de informes de cumplimiento y auditorías internas.
\item Automaticé tareas de monitorización con scripts en Python y Bash.
\end{itemize}
\textbf{Palabras clave}: Análisis de riesgos, Cumplimiento normativo, Cortafuegos, IDS/IPS, Python}
\section{Idiomas}

\cvline{Español}{Competencia nativa o bilingüe}
\cvline{Inglés}{Competencia profesional plena \hfill \textbf{Palabras clave}: C1, TOEFL 110}
\cvline{Francés}{Competencia limitada de trabajo \hfill \textbf{Palabras clave}: A2}
\section{Competencias}

\cvline{Seguridad Ofensiva}{Experto \hfill \textbf{Palabras clave}: Pentesting, OSCP, Metasploit, Burp Suite, Nmap}
\cvline{Seguridad Defensiva}{Avanzado \hfill \textbf{Palabras clave}: SIEM, Splunk, Wireshark, IDS/IPS, Cortafuegos}
\cvline{Programación}{Avanzado \hfill \textbf{Palabras clave}: Python, Bash, PowerShell, SQL}
\cvline{Seguridad en la Nube}{Intermedio \hfill \textbf{Palabras clave}: AWS, Azure, IAM, Kubernetes}
\section{Premios}

\cventry{nov 2021}
        {INCIBE}
        {Premio Nacional de Ciberseguridad}
        {}
        {}
        {Reconocimiento por la identificación de una vulnerabilidad crítica en una infraestructura pública.}

\cventry{jul 2019}
        {Universidad Politécnica de Madrid}
        {Mejor Proyecto Fin de Grado}
        {}
        {}
        {Proyecto destacado en detección de intrusiones mediante machine learning.}
\section{Certificados}

\cventry{mar 2023}
        {ISC²}
        {CISSP}
        {\url{https://www.isc2.org/certifications/cissp}}
        {}
        {}

\cventry{jun 2021}
        {Offensive Security}
        {OSCP}
        {\url{https://www.offensive-security.com/pwk-oscp/}}
        {}
        {}

\cventry{sept 2020}
        {EC-Council}
        {CEH}
        {\url{https://www.eccouncil.org/programs/certified-ethical-hacker-ceh/}}
        {}
        {}
\section{Publicaciones}

\cventry{ene 2022}
        {Análisis de vulnerabilidades en infraestructuras críticas}
        {Revista Española de Ciberseguridad}
        {\url{https://www.revistaciberseguridad.es/articulo}}
        {}
        {Artículo sobre metodologías de pentesting en entornos industriales y la importancia de la segmentación de redes.}
\section{Proyectos}

\cventry{feb 2021–ago 2021}
        {Desarrollo de un escáner de vulnerabilidades para aplicaciones web}
        {Herramienta de análisis de vulnerabilidades web}
        {\url{https://github.com/cmendoza-sec/vuln-scanner}}
        {}
        {\begin{itemize}
\item Escáner automatizado para detectar fallos de seguridad en aplicaciones web.
\item Implementación de pruebas basadas en OWASP Top 10.
\item Interfaz de línea de comandos y generación de informes en formato PDF.
\end{itemize}
\textbf{Palabras clave}: Python, OWASP, Seguridad Web, Automatización}
\section{Intereses}

\cvline{Deportes}{Ciclismo, Senderismo, Natación}
\cvline{Música}{Guitarra, Piano}
\section{Voluntariado}

\cventry{sept 2020 hasta la fecha}
        {Mentor en Ciberseguridad}
        {Fundación Ciberseguridad para Todos}
        {\url{https://www.ciberseguridadparatodos.org}}
        {}
        {\begin{itemize}
\item Mentoría a jóvenes en riesgo de exclusión social en fundamentos de seguridad informática.
\item Organización de talleres prácticos sobre contraseñas seguras y navegación segura.
\item Fomento de vocaciones STEM entre estudiantes de secundaria.
\end{itemize}}
    
\end{document}